\documentclass[11pt]{article}
\usepackage[margin=1.15in]{geometry}
\usepackage{amsmath,amssymb,amsthm}
\usepackage{booktabs}
\usepackage[colorlinks=true,linkcolor=black,citecolor=black,urlcolor=black]{hyperref}
\setlength{\parskip}{0.55em}
\setlength{\parindent}{0pt}
\newtheorem{proposition}{Proposition}
\theoremstyle{definition}
\newtheorem{definition}{Definition}
\theoremstyle{remark}
\newtheorem{remark}{Remark}
\newcommand{\Ex}{\mathbb{E}}
\renewcommand{\Pr}{\mathbb{P}}

\title{Inheritance of Possibility and Demonstrated Function}
\author{Flyxion\\Independent Researcher}
\date{October 2026}

\begin{document}
\maketitle

\begin{abstract}
Ecological and resilience research often establishes that a system has an enabling condition, such as an inherited substrate or a constructed capacity, and then describes the system as performing the function that condition would permit. The two statements differ in evidential content. This essay defines the difference: enabling conditions relative to a background, realized function stated so that it could have failed, legacy-mediated function as distinct from functional inheritance, and a scheme of claim types with a categorical claim gap. Seven short propositions show where common designs leave the stronger claims unidentified. Mineral evolution serves as a reference standard for stating an enabling-condition framework without asserting more than it shows. A study of submerged peat substrates and a study of an urban resilience index serve as worked cases. The conclusion is deliberately narrow: the observed substrates and constructed capacities do not by themselves identify the proposed functions, and whether either study demonstrates enabling effects, realized performance, or causal inheritance must be settled claim by claim.
\end{abstract}

\section{Introduction}

A large class of empirical papers has the following shape. A condition is measured. The condition is one that would make some process possible, or more probable, or more intense. The paper then speaks in the vocabulary of the process: the substrate \emph{supports} a community, the habitat \emph{serves as} a nursery, the capacity \emph{confers} resilience. The measurement is real and often careful. The vocabulary, however, belongs to a different class of claim, one that requires evidence about rates, interactions, counterfactuals, or independent outcomes that the measurement of the condition does not supply.

This essay does not argue that such claims are false. It argues that they occupy a different evidential position from the observations offered for them, and that the difference can be stated precisely enough to be checked. The vocabulary of ``possibility'' and ``function'' is used loosely in the literature, so the first task is definitional. The second is to show, by short formal arguments, which inferences a given design can and cannot support. The third is to apply both to a reference standard and to two cases. Mineral evolution is used as the reference standard because it states its terms and the relation between an enabling history and an observed set explicitly enough that its revisions can be compared, which is what the two cases lack. The treatment is general: the cases are examples of a structure, and nothing is claimed about the field as a whole beyond what the two examples and the cited literature support.

\section{Definitions}

\subsection{Conditions, backgrounds, and functions}

Let a system be described by a state $\sigma$ in a state space $\Sigma$, and let $z$ denote a background, the collection of other variables (colonization sources, disturbance regime, institutional setting) that are held fixed when a condition is considered. A \emph{condition} is a binary or graded property $c$ of the state. A \emph{process} $p$ is a family of trajectories in $\Sigma$, and a \emph{function} is a process together with an output $\phi$ defined on it and measured in the units of the process (for example, individuals recruited per unit area per year, or productivity retained after a disturbance). Let $A(c,z)$ be the set of processes that are physically admissible given the condition and background, and let $R(a,z)\subseteq A(a,z)$ be the processes realized when an intervention assigns $c=a$ under background $z$.

\begin{definition}[Enabling condition]\label{def:enabling}
Relative to a background $z$, the condition $c$ is \emph{strongly enabling} for a process $p$ if $A(c{=}0,z)\subsetneq A(c{=}1,z)$ and $p\in A(c{=}1,z)\setminus A(c{=}0,z)$. It is \emph{weakly enabling} if instead $\Pr\bigl[p\in R(1,z)\bigr]>\Pr\bigl[p\in R(0,z)\bigr]$, the probability being taken over the unspecified further variables. Both are interventional contrasts; an observational association between $c$ and realization identifies them only under additional assumptions, such as the absence of confounding of $c$ given $z$.
\end{definition}

The distinction matters because a substrate may make recruitment more likely without making an otherwise impossible process possible, and many ecological conditions enable only in conjunction with others, which is why the definition is relative to $z$ and compares matched states differing in $c$. Neither form places a process in the realized set. Realization depends on further variables that the condition does not fix.

\begin{definition}[Realized function]\label{def:realized}
A function $\phi$ is \emph{realized} in a state if the process $p$ is in $R$ and $\phi$ has been measured, with its uncertainty, by a procedure that is observed independently of the rule that assigns the condition or capacity score, so that the observed value is not algebraically entailed by that score.
\end{definition}

This definition does not forbid a measurement from using the condition in ordinary ways. Recruitment per unit substrate area must use substrate area, and an outcome estimator may adjust for a measured capacity. What it excludes is definitional dependence, where the purported outcome is assigned by the same rule that constructs the condition, so that agreement is a property of the construction. Realization requires an observation that could have disagreed. It follows from Definitions~\ref{def:enabling} and \ref{def:realized} together that observing a realized function on a substrate does not, by itself, establish that the substrate is enabling, because enabling is an interventional contrast and the observation supplies no contrast; this is exactly the inference the essay examines.

\subsection{Histories, continuity, and inheritance}

Write $\sigma_1(h)$ for the later state generated under history $h$, with $h$ the actual history, which includes an earlier system $\sigma_0$, and $h^{\varnothing}$ a counterfactual history from which a named element of $\sigma_0$ has been removed. If $\sigma_1$ is a complete state description, two histories ending in the same state cannot differ in admissible processes unless history belongs to the description, so $\sigma_1$ is here understood as a partial observation of a site, and the counterfactual comparison is a comparison of potential outcomes in the sense of Rubin \cite{rubin1974}. The substantive question is what the counterfactual removes and what it preserves. Removing peat while preserving a consolidated structure of the same form tests the distinctive effect of peat. Removing the inherited structure altogether tests the total effect of the material legacy. Neither estimand should silently stand for the other.

\begin{definition}[Material continuity]\label{def:material}
Material continuity is the fraction of a later material pool that derives from an earlier one. A source-partitioning calculation estimates it only under an added attribution assumption. As an example, for a two-end-member mixing model with tracer values $\delta_s$ (sample), $\delta_{\mathrm{marine}}$ and $\delta_{\mathrm{terr}}$,
\[
f_{\mathrm{terr}}=\frac{\delta_s-\delta_{\mathrm{marine}}}{\delta_{\mathrm{terr}}-\delta_{\mathrm{marine}}},
\]
which estimates the share of material from terrestrial rather than marine sources under the model. It measures historical continuity only when the terrestrial end member can independently be attributed to the earlier pool, since later terrestrial inputs could give the same signature. It is meaningful only when the end members are separated relative to measurement error, and with more than two unconstrained sources a single tracer generally does not uniquely identify the fractions \cite{parnell2010}.
\end{definition}

\begin{definition}[Functional continuity]\label{def:functional}
Functional continuity between an earlier state $\sigma_0$ and a later state $\sigma_1$ holds for a function $\phi$ if $|\phi(\sigma_1)-\phi(\sigma_0)|\le\epsilon$, where the tolerance $\epsilon$ is fixed in advance from measurement error or an equivalence margin and not from the observed difference, and the two values rest on the same process, measured in the same units and by comparable methods.
\end{definition}

Fixing $\epsilon$ in advance is necessary because for any two finite values some tolerance exceeds their difference, so an existentially quantified tolerance asserts nothing.

\begin{definition}[Inheritance of possibility]\label{def:inhpos}
The later state \emph{inherits possibility} from the earlier system if, holding the background fixed, $A\bigl(\sigma_1(h)\bigr)\supsetneq A\bigl(\sigma_1(h^{\varnothing})\bigr)$.
\end{definition}

\begin{definition}[Legacy-mediated function and functional inheritance]\label{def:funcinh}
A function $\phi$ is \emph{legacy-mediated} in $\sigma_1$ if (a) $\phi$ is realized in $\sigma_1(h)$ and (b) $\phi(\sigma_1(h))-\phi(\sigma_1(h^{\varnothing}))>0$ for the stated counterfactual. It is \emph{functionally inherited} from $\sigma_0$ if, in addition, (c) $\phi$ was realized in $\sigma_0$ and (d) the mechanism $\mu$ by which $\phi$ is produced in $\sigma_1$ corresponds to the mechanism by which it was produced in $\sigma_0$, in the sense that the two operate on the same causal pathway.
\end{definition}

One illustrative criterion for (d) is correspondence of mediators: let $M_0$ and $M_1$ be the sets of variables through which the function is carried in $\sigma_0$ and $\sigma_1$ (for example, a detrital food-web pathway fed by organic carbon, as against provision of consolidated hard structure for settlement). Here $M_1$ matches $M_0$ when the same mediating variables, identified by the same measurement, carry the effect in both states; the weaker reading, a similar causal graph up to relabeling, is not enough. Condition (d) holds to the extent that $M_1$ matches $M_0$ in this sense and that intervening on a shared mediator changes $\phi$ in both states in the same direction. In the peat case, a function carried by the hard structure alone would be shared with any consolidated substrate and would not correspond to a wetland's detrital function, whereas one carried by the organic-carbon pathway might. Condition (d) is needed because an old forest could support biodiversity, leave a physical remnant, and thereby support a wholly different community through a wholly different mechanism. That satisfies (a) to (c) if $\phi$ is broad enough, and it may reasonably count as a legacy-mediated function, but functional inheritance, as opposed to a legacy effect, requires a relation between the earlier and later mechanisms. In this vocabulary, persistence of the original ecosystem's function is functional continuity in the sense of Definition~\ref{def:functional}, whereas occupancy of inherited material by a new community establishes neither legacy mediation nor functional inheritance by itself, since the community might or might not preserve the relevant causal pathway and that has to be shown. A non-peat hard-substrate control, matched on the covariates that govern colonization, estimates the distinctive-effect counterfactual in (b). Even then, demonstrating that an \emph{earlier} function was inherited requires historical evidence about $\sigma_0$ for (c) and (d), because a control comparison speaks only to the present substrate effect. Niche-construction theory and the review literature on engineering legacies give the general setting for such questions \cite{laland2016,albertson2024}.

\subsection{Claim types and the claim gap}

Claims in this literature differ in what must be shown for them to be true. They can be arranged by the kind of evidence they require, though not as a single entailment chain: occupancy does not require persistence of the original system, interaction evidence does not necessarily precede a demonstrated export flow, and an urban disturbance outcome has no natural counterpart to ecological occupancy. The scheme is therefore not a chain of entailments but a classification of claim types by the kind of evidence each requires, in two branches, shown in Table~\ref{tab:ladder}. The order within a branch is a convention of increasing evidential demand and not a theorem: a marked export flow, for instance, could establish $C_4$ without observing the interactions that produced it, so $C_4$ need not rest on $C_3$.

\begin{table}[h]
\centering
\small
\begin{tabular}{@{}p{0.07\textwidth}p{0.34\textwidth}p{0.47\textwidth}@{}}
\toprule
Type & Claim & Evidence required \\
\midrule
$C_1$ & A substrate or structure persists or has been built & Direct characterization (dating, mapping, sourcing of materials) \\
$C_2$ & Organisms occupy it & Occupancy or abundance data with adequate sampling \\
$C_3$ & The occupants interact, so the assemblage is an ecosystem in the working sense & Interaction data (trophic, behavioural, demographic) \\
$C_4$ & The system contributes to other habitats (functional extension) & Rate or flow measurements: movement, survival, recruitment \\
$C_5$ & The contribution depends on the inherited feature & Matched controls and, for inheritance of an earlier function, historical evidence \\
\midrule
$U_1$ & A set of capacities is defined and scored & Indicator construction and aggregation rules \\
$U_2$ & The score changes over time or differs among units & Repeated measurement with an uncertainty statement \\
$U_3$ & The score is associated with an independent outcome & Outcome data not constructed from the indicators \\
$U_4$ & The score predicts independent outcomes and performs favourably against specified alternative aggregations and component-based models & Multiple independent events, with evaluation on held-out events \\
\bottomrule
\end{tabular}
\caption{Claim types. $C_i$ is the ecological branch and $U_i$ the urban resilience branch. The order within a branch is by evidence required, and cross-branch correspondence is by analogy of evidential structure, not shared vocabulary: $C_1\leftrightarrow U_1$, $C_4\leftrightarrow U_3$, $C_5\leftrightarrow U_4$.}\label{tab:ladder}
\end{table}

The correspondences are structural. $C_1$ and $U_1$ both concern characterization of the enabling condition or capacity itself. $C_4$ and $U_3$ both require a quantity about something other than the system's own structure, measured without reference to how the condition was characterized. $C_5$ and $U_4$ both require that the result depend on the specific feature or form claimed. A paper that spans both branches, such as an urban ecological resilience study, is coded claim by claim, each claim against the branch that matches its subject: claims about habitat structure and occupancy against the $C$ types, and claims about an aggregated capacity score against the $U$ types.

\begin{definition}[Claim gap]
For a paper, let $S_{\mathrm{claimed}}$ be the set of claim types asserted in its title, abstract, or discussion and $S_{\mathrm{supp}}$ the set whose evidence requirements are met by the reported analysis. The \emph{claim gap} is the set $G=S_{\mathrm{claimed}}\setminus S_{\mathrm{supp}}$.
\end{definition}

The gap is reported as a set, because its size has no interval meaning: a gap of two types is not demonstrably twice the evidential deficiency of a gap of one, and a scalar such as $|G|$ or a difference of levels is only a coding convenience. A nonempty $G$ is not an error in any measurement; it is a property of the relation between measurements and sentences. The set $S_{\mathrm{claimed}}$ is read from the text, and $S_{\mathrm{supp}}$ is judged by comparing the reported design and data with the evidence requirements in Table~\ref{tab:ladder}; the propositions below are tools for that second judgment and do not by themselves compute $G$. Because $G$ is defined by a coding rule applied to text, two coders could disagree, and any empirical use of it would need to report agreement between coders.

As an illustration of coding, take the peat abstract's sentence that the substrate provides structure ``for fish communities and as nursery habitat''. The first clause codes as $C_2$, since fish occupy the structure, and the second as $C_4$, since a nursery role is a contribution to other habitats. The surveys reach $C_1$ and $C_2$, so for this sentence $S_{\mathrm{claimed}}=\{C_2,C_4\}$, $S_{\mathrm{supp}}=\{C_1,C_2\}$ for the paper as a whole, and the gap is $G=\{C_4\}$. The abstract's reference to a detritus-based food web would add $C_3$ to the claimed set. The coding is a judgment that a second coder could dispute, which is why agreement between coders would have to be reported.

\section{What the Designs Can and Cannot Identify}

The following propositions are elementary. Their value lies in making explicit what a given design leaves open. They are tools for distinguishing identification limits (Propositions~\ref{prop:overlap} to \ref{prop:support}), sampling resolution (Propositions~\ref{prop:perm} and \ref{prop:kish}), and consequences of index construction (Propositions~\ref{prop:elast} and \ref{prop:decomp}). The first group states that a class of data cannot distinguish a state with $c{=}1$ from one with $c{=}0$ in its effect on $R(a,z)$ or on the function $\phi$, or cannot recover a quantity on which a claim depends. The second concerns what a test can resolve. The third concerns an algebraic property and an accounting identity of the index.

\begin{proposition}[Overlap does not identify movement]\label{prop:overlap}
Let two assemblages of $n_a$ and $n_b$ species be drawn independently and uniformly from a common filter set of $f$ species, with $\max(n_a,n_b)\le f\le K$, where $K$ is the regional pool. The overlap $O$ is hypergeometric,
\[
\Pr(O=o\mid f)=\frac{\binom{n_a}{o}\binom{f-n_a}{n_b-o}}{\binom{f}{n_b}},\qquad \Ex[O]=\frac{n_an_b}{f},
\]
and $\Ex[O]$ ranges from $n_an_b/K$ to $\min(n_a,n_b)$ as $f$ varies. The support of $O$ is $\max(0,n_a+n_b-f)\le o\le\min(n_a,n_b)$. Every overlap within this support has positive probability under the no-movement model. For $O>0$, its expectation is matched when $n_an_b/O$ is an admissible integer pool size, and is approximated by a neighbouring admissible integer when $n_an_b/O$ lies within the permitted range $[\max(n_a,n_b),K]$. If it exceeds $K$, rounding cannot repair the incompatibility, and the no-movement model with that pool cannot match the expectation.
\end{proposition}

\emph{Argument.} Each species of the first assemblage is in the second with probability $n_b/f$, so the expectation follows by linearity, and the mass function is the standard hypergeometric count. \qed

For small-fish assemblages of 55 and 59 species with 41 shared (the partition reported for the peat study discussed below), the expectation is matched at $f=n_an_b/O\approx79.1$, which is not an integer. At $f=79$ the expected overlap is $41.08$, the probability of exactly 41 shared species is $0.22$, and the probability of at least 41 is $0.62$. The observed overlap is therefore an ordinary outcome of two draws from a common pool of about 79 species with no connectivity process, where that pool is an illustrative null model and not an independently established ecological quantity. The common-pool null is one of a family of no-movement nulls. Draws weighted by an abundance distribution, for instance, give different expected overlaps, and the proposition's point is only that some null in the family can match the observed overlap, so overlap alone does not discriminate between them and movement. The proposition does not show that connectivity is absent. Connectivity and common filtering (depth range, region, sampling method) can each generate overlap, and the defensible conclusion is that shared composition alone does not identify movement. The same reasoning applies when the shared set is taken from a checklist of target species.

\begin{proposition}[Density does not identify nursery contribution]\label{prop:nursery}
In a simplified rendering of the contribution framework of Beck et al.\ \cite{beck2001}, the contribution of habitat $h$ to the adult stock is $c_h=A_h\,\rho_h\,s_h\,\tau_h$, where $A_h$ is habitat area, $\rho_h$ is juvenile density (individuals per unit area), $s_h$ is survival to recruitment (dimensionless), and $\tau_h$ is the fraction of survivors reaching adult habitat. Per unit area, the contribution is $\rho_hs_h\tau_h$. Density measurements identify $\rho_h$ only: for any ordering $\rho_1>\rho_2$ and any ordering of per-area contributions there exist $(s_h,\tau_h)$ compatible with it.
\end{proposition}

\emph{Argument.} To obtain $c_1<c_2$ per unit area with $\rho_1>\rho_2$, choose $s_1\tau_1<(\rho_2/\rho_1)s_2\tau_2$. Nothing in density data constrains $s_h$ or $\tau_h$. \qed

The framework itself defines nursery value comparatively, as greater production of recruits per unit area than other habitats, and the critical literature on nursery inference has stressed that presence of juveniles is not equivalent to contribution \cite{sheaves2006,dahlgren2006}. The proposition states that point in the units that carry it.

\begin{proposition}[Substrate effects without common support]\label{prop:support}
Let the response satisfy $Y=\alpha+\beta S+g(Z)+\varepsilon$ with substrate $S\in\{0,1\}$ and covariate $Z$ (for example depth), and suppose the supports of $Z$ in the two substrate groups are disjoint. Then $\beta$ is identified only by an assumption about the form of $g$ outside the observed support of each group.
\end{proposition}

\emph{Argument.} On disjoint supports, any value of $\beta$ can be offset by a modification of $g$ on one support, leaving the observed conditional means unchanged. \qed

This is a violation of the positivity condition familiar from causal inference \cite{petersen2012}. Its relevance is limited: it concerns the substrate coefficient and not the existence of the assemblage.

\begin{proposition}[Permutation floor]\label{prop:perm}
For $N$ units divided into groups of sizes $n_1,\dots,n_J$, the number of distinct assignments is $M=N!/\prod_jn_j!$, so an exact one-sided permutation test cannot return a $p$-value smaller than $1/M$ when the observed assignment is the unique most extreme one. For two groups of equal size $k$ and an absolute-difference statistic, complementary assignments are equally extreme, so the two-sided floor is $2/\binom{2k}{k}$.
\end{proposition}

For groups of size $(1,3)$, $M=4$ and the one-sided floor is $0.25$. For groups $(5,3)$, $M=56$ and the one-sided floor is $1/56\approx0.018$. For $k=3$ the two-sided floor is $0.1$, and for $k=4$ it is $2/70\approx0.029$. These are statements about attainable resolution under a stated exchangeability scheme, and not about power, confounding, or the precision needed to resolve an effect of biological interest, so they are not a universal requirement on the number of sites.

\begin{proposition}[Design effect for nested replicates]\label{prop:kish}
Suppose $m$ replicates within a unit have a common variance and an exchangeable intraclass correlation $\rho$. The variance-equivalent sample size for estimating the unit mean is $n_{\mathrm{eff}}=m/[1+(m-1)\rho]$, which tends to $1$ as $\rho\to1$.
\end{proposition}

This is the Kish design effect, $1+\rho(m-1)$ for clusters of size $m$ \cite[Sec.~5.4]{kish1965}. It is a statement about estimating a mean and not a licence to treat repeated surveys at one site as independent replication of a substrate treatment, a distinction that goes back to Hurlbert \cite{hurlbert1984}. Whether the concern bites in a particular analysis depends on the inferential target, a point argued by Oksanen \cite{oksanen2001} and by Davies and Gray \cite{davies2015}, and $\rho$ is an empirical quantity to be estimated or bounded and not assumed.

\begin{proposition}[Elasticity in a multiplicative index]\label{prop:elast}
Let $I=D\,(B+E+T+BET)$ with $B,E,T,D>0$, and $\varepsilon_X=\partial\ln I/\partial\ln X$. Then $\varepsilon_D=1$ and
\[
\varepsilon_B=\frac{B(1+ET)}{B+E+T+BET}<1,
\]
and symmetrically for $E$ and $T$. The inequality is strict on the positive domain, with equality only in the excluded limit $E+T=0$.
\end{proposition}

\emph{Argument.} Differentiate $\ln I=\ln D+\ln(B+E+T+BET)$; the bound follows from $B(1+ET)=B+BET<B+E+T+BET$ when $E+T>0$. \qed

The proposition says that the specification hard-codes the ordering of proportional elasticities, with $D$ largest, irrespective of data. The property is not specific to this index or a criticism of the authors' choice: any index of the form $D\,g(B,E,T)$ has $\varepsilon_D=1$ identically, and the other elasticities fall below one whenever those of $g$ do, which the authors acknowledge by calling the specification a modelling assumption. The proposition does not say that $D$ is a binding constraint, because the largest sensitivity to a proportional change, the largest contribution to observed variation, a binding constraint in the sense of a limiting factor, and the most beneficial intervention are different quantities.

\begin{proposition}[Decomposition of an index change]\label{prop:decomp}
For units $i=1,\dots,n$ let $I_i=D_iP_i$ with $P_i=B_i+E_i+T_i+B_iE_iT_i$, and use the population covariance (divisor $n$). Then $\bar I=\bar D\bar P+\operatorname{Cov}(D_i,P_i)$, and hence $\Delta\bar I=\Delta(\bar D\bar P)+\Delta\operatorname{Cov}(D_i,P_i)$.
\end{proposition}

This is an accounting identity and not a causal explanation. It separates the part of a change that comes from the product of mean capacities from the part that comes from their covariance across units.

\section{A Reference Standard: Mineral Evolution}

The programme of mineral evolution \cite{hazen2008,hazen2010} classifies mineral species by the paragenetic modes through which they form and traces the diversification of those modes through planetary history. It is useful as a standard for three reasons. It states the relation between an enabling history and an observed set explicitly enough to be checked. It makes its definitions explicit enough that revisions can be compared. And it distinguishes in practice between what is observed (species and modes in the record) and what is inferred about origin.

The standard also illustrates the thesis about ambiguous terms. The 2008 review, working with about 4,300 known species, stated that biological activity, chiefly through the surface oxygenation it brought about, ``may be responsible directly or indirectly for more than two thirds of all known mineral species'', and its abstract put the point as ``most''; the wording is hedged and the mechanism is itself an enabling condition, since oxygenation made weathering products possible without being the sole cause of any one species. The 2022 study of Hazen and Morrison, which surveyed 57 paragenetic modes among 5,659 mineral species, reports that at least 2,707 species (47.8 percent) form as a result of biological influences, mostly oxidative weathering, and that only about 34 percent, more than 1,900 species, form exclusively as a consequence of biological processes. It states that this exclusive figure is contrary to previous estimates, ``in contrast to the estimate of two-thirds by Hazen et al.\ (2008)'' \cite{hazen2022}. The quantities compared are different. Let $\mathcal{B}$ be the set of bio-mediated modes and $M(s)$ the set of modes by which species $s$ forms, and define
\[
B_{\mathrm{any}}=\{s: M(s)\cap\mathcal{B}\neq\varnothing\},\qquad B_{\mathrm{ex}}=\{s: M(s)\subseteq\mathcal{B}\}.
\]
Then $B_{\mathrm{ex}}\subseteq B_{\mathrm{any}}$. The 2008 statement concerns something close to $B_{\mathrm{any}}$ (``directly or indirectly''), the 2022 exclusive figure is $|B_{\mathrm{ex}}|/N$, and the 2022 figure of 47.8 percent is the nearer counterpart to the 2008 one, though the two also differ in the species inventory (about 4,300 against 5,659) and in the treatment of anthropogenic and taphonomic minerals. A figure for one set cannot be compared with a figure for the other without matched definitions, and the 2022 paper's own juxtaposition of the exclusive figure with the 2008 estimate is an instance of the ambiguity. The quoted statements are taken from the text of each paper, so the comparison is a reading of what each says and not a finding about the earlier estimate's correctness. The point for present purposes is that the framework states its terms, and so permits a reader to see exactly when two statements are comparable.

\section{Case A: Submerged Peat as Reef Habitat}

\subsection{Enabling condition and claim}

Lebrato and coauthors \cite{lebrato2026} report submerged tropical peat reefs in the Bazaruto Archipelago, radiocarbon dated to about 300 to 5,000 years before present, with organic carbon up to 40 percent by dry weight and stable-carbon values consistent with a terrestrial plant origin. The title states that ancient terrestrial forests are ``reborn as modern marine ecosystems''. The abstract states that the peat substrate provides structure ``for fish communities and as nursery habitat'', that overlapping communities ``highlight ecological connectivity across reef types'', that the system hosts ``a detritus-based food web'', and that ancient terrestrial systems ``can persist and support contemporary ecological functions''. The Discussion is more guarded on the nursery point: the authors encountered numerous juveniles on peat reefs, which ``suggests a nursery function''. In the vocabulary of Section~2, the substrate characterization reaches $C_1$, the surveys reach $C_2$, and the title and abstract invoke $C_3$ through $C_5$.

\subsection{Where identification stops}

The points below are of two kinds, and they should not be conflated. The overlap, trophic, nursery, and substrate-effect points are identification limits: the data collected cannot distinguish the claim from rivals however many more of the same observations were made. The replication and permutation points concern resolution and reporting: they limit how sharply a test can speak, and they are matters for clarification or for added sampling and not fatal design flaws.

\emph{Overlap.} The small-fish Venn partition gives 14 coral-only, 18 peat-only, and 41 shared species, so 55 coral species and 59 peat species pooled across sites. The Results text gives 54 coral species, a one-species discrepancy with the partition. By Proposition~\ref{prop:overlap} the overlap is an ordinary outcome of a common pool of about 78 to 79 species. The peat figure is pooled over several reefs, with 26 to 36 species at individual reefs, whereas the coral figure comes from a single site, so unequal sampling effort also enters.

\emph{Megafauna overlap.} The megafauna partition gives no coral-only species, one peat-only species, and 44 shared, which implies 44 species at the coral site, yet the text reports that coral reef surveys recorded 25 species. The surveys targeted 46 pre-determined indicator species, and the displayed sets may be drawn from that list, though the paper does not say so. Either way the displayed shared set cannot be read as a count of species co-occurring at the surveyed sites. The authors themselves note that the multivariate tests for megafauna were not significant despite $R=1.000$ and attribute this to limited sampling.

\emph{Trophic claim.} The abstract refers to a detritus-based food web, a claim of type $C_3$. The survey methods described are benthic quadrats, invertebrate belt transects, small-fish visual census, and megafauna transects, together with peat geochemistry, and I find among them no measurement of trophic interactions, so the food-web statement is an inference from composition and substrate and not an observation.

\emph{Nursery role.} The evidence offered is the presence of juveniles. Presence establishes occurrence, and adequately sampled counts could estimate density. The reported small-fish abundance, including the smallest size classes of under 5 cm and 5 to 10 cm, was higher at the coral site than on the peat reefs. These counts do not supply a comparative juvenile-density advantage for peat reefs, though they do not contradict a nursery contribution either, since size classes need not identify juveniles of the same species, density is only one component of contribution, and the comparison is of one coral site against several peat reefs. By Proposition~\ref{prop:nursery}, density at most identifies $\rho_h$, whereas survival, growth, and transfer to adult habitat carry the contribution.

\emph{Substrate effect.} The peat reefs lie at 2 to 4 m depth and the coral reference site at 12 to 22 m in the Results (11 to 16 m in a figure caption, an internal inconsistency). The depth ranges do not overlap, so by Proposition~\ref{prop:support} differences in assemblage cannot be assigned to substrate rather than depth without an assumption about the depth effect. The comparison is also thinly replicated. The coral habitat is a single reference site surveyed on five occasions, and peat reefs received one survey each, so site-level comparisons have groups of size $(1,3)$ with floor $0.25$ under Proposition~\ref{prop:perm}, and survey-level comparisons $(5,3)$ with floor $0.018$. The Results report an ANOSIM $p=0.013$ for invertebrate species composition. As a diagnostic only: if the eight surveys were the exchangeable units and the probability were calculated by exhaustive enumeration of the 56 assignments, the reported value would fall below the exact floor of $1/56$. Monte Carlo permutation sampling or another calculation convention could produce an estimate below that floor, and the paper does not state the permutation units, the admissible assignments, or the probability calculation, so the value calls for clarification and not a finding of error. The benthic quadrat result ($n=134$ quadrats, $p=0.001$) has quadrats as its units of observation, and unless the permuted units are documented the analysis is informative about quadrat-scale variation within the sampled reefs but not replication of reef type, with the qualification of Oksanen \cite{oksanen2001} and Davies and Gray \cite{davies2015} that the inferential target determines whether this matters.

\subsection{Functional inheritance in this case}

The peat case is the one in which material inheritance is literal: the substrate is the remnant of an earlier system. Material continuity in the sense of Definition~\ref{def:material} could in principle be quantified by a mixing fraction, and the persistence of a consolidated structure is a $C_1$ claim. Definition~\ref{def:funcinh} requires more. For a legacy-mediated function, a realized function and a positive effect of the history relative to a stated counterfactual are needed, and for functional inheritance, evidence about the earlier system's function and about correspondence of mechanisms. The study documents an assemblage on inherited material. Whether the peat contributes something beyond consolidated structure is the question a hard-substrate control would address, and whether an \emph{earlier} function persisted would need additional historical evidence. A new community performing a function by occupying inherited material is a coherent and interesting finding. It is not yet functional inheritance, and the abstract's statement that ancient terrestrial systems persist and support contemporary functions is a claim of a stronger kind than the surveys support.

\section{Case B: Constructed Potential and Independently Demonstrated Performance}

\subsection{Setting}

Hou, Hu and Halik \cite{hou2026} construct an urban resilience index, TE4, for ten city and county units on the northern slope of the Tianshan Mountains over 2014 to 2024. It combines resistance $B$, adaptive capacity $E$, institutional learning and preparedness $T$, and recovery $D$ as $\mathrm{TE4}_i=D_i(B_i+E_i+T_i+B_iE_iT_i)$. The paper was accepted on 7 October 2026 and appears as an Article in Press, so the text discussed here may change before final publication. The case differs from Case A in that the enabling condition is constructed and not inherited, so it does not concern literal inheritance of a substrate; it concerns the relation between a constructed potential, which the index measures, and performance in a disturbance, which it is meant to characterize. It shares with Case A the structure that an enabling condition has been characterized and the question is whether it yields a demonstrated function.

\subsection{What the study establishes and what it states}

The study is unusually candid. It describes the multiplicative specification as a theoretical modelling assumption, notes that recovery is the principal modelled deficit ``by design'' and not a universal causal gatekeeper, treats its external check as exploratory, and describes TE4 as a structural capacity diagnostic and not a predictor of annual productivity. Those hedges are appropriate to what it shows, and the claim gap for this paper is accordingly small.

Three points follow from the formal tools. First, because $D$ enters multiplicatively, Proposition~\ref{prop:elast} implies that $D$ has the largest proportional elasticity irrespective of data, which the authors acknowledge, and so the identification of recovery as the principal deficit is partly a consequence of the form and says nothing yet about which capacity would most improve outcomes. Second, the reported endpoint change decomposes by Proposition~\ref{prop:decomp}. Using the reported values, TE4 fell from $0.1184$ to $0.1168$ (a change of $-0.0016$), mean recovery from $0.3723$ to $0.3424$, and potential resilience rose from $0.3149$ to $0.3441$. Taking the potential values as cross-unit means of $P_i$, the product of means changes by $+0.00058$ and the covariance term by $-0.00218$. Allowing for rounding of all six inputs to four decimals, the first lies between about $+0.00051$ and $+0.00065$ and the second between about $-0.00235$ and $-0.00201$, so the signs are robust (the calculation is in the Appendix). If the potential values were instead computed from mean capacities, the residual is not a covariance in the strict sense, but the accounting remains valid as a residual. The paper itself describes these as endpoint changes and not statistical trends, and reports a non-significant trend in TE4, so the decomposition should not be read as a temporal finding. Third, the independent check, annual productivity under drought, is the step from $U_1$ toward $U_3$. The paper reports that higher TE4 did not attenuate short-term productivity losses under stronger drought, and describes this as ``confirming its role as a structural capacity-diagnostic index''; the event-level association between pre-event recovery capacity and post-drought recovery rests on ten valid events, is threshold-sensitive, and is not statistically conclusive under the primary threshold. A null association is compatible with the stated conceptual separation between structural capacity and short-term productivity, and equally compatible with an index that has no predictive content for that outcome, so the null does not by itself discriminate between them.

\subsection{Constructed capacities are not validated possibilities}

A capacity index is a $U_1$ object. Reviews of resilience measurement have repeatedly noted that agreement of indices with independent outcomes is rarely demonstrated. Carpenter and coauthors framed the problem as specifying resilience of what to what \cite{carpenter2001}; Bakkensen and coauthors noted that few resilience and vulnerability indices had been tested against observed disaster impacts and found, for five leading U.S. indices, that four explained damages as theory predicts, three explained fatalities, and two explained disaster declarations \cite{bakkensen2017}; Gu and coauthors tested indicators and frameworks against post-disaster outcomes for 1,716 U.S. counties and found that complex or hierarchical designs did not reliably represent resilience better than simpler ones \cite{gu2023,gu2025}; and Halekotte, Mentges and Lichte argue that capacity assessments should rely on multiple performance curves and validation against real or simulated events \cite{halekotte2024}. A constructed capacity should not be treated as an already validated possibility. Because the index is a deterministic function of its components, it cannot carry information beyond their joint values. What it can show is useful compression, a defensible functional form, or better prediction than specified alternatives. A constructed capacity becomes a validated possibility when it predicts independent outcomes on held-out events and performs favourably against specified alternative aggregations and component-based models.

The positive side can be stated concretely. Useful compression would mean that one scalar retains most of the predictive information of the four capacities: on held-out disturbance events, TE4 predicts post-disturbance outcomes about as well as a model that enters $B$, $E$, $T$, and $D$ separately, with fewer parameters, and better than a linear aggregate and than the symmetric alternative the authors define. A defensible functional form would mean that the multiplicative recovery term improves out-of-sample prediction over versions that do not single out $D$, so that the asymmetry is supported by outcomes and not only by the theoretical rationale the authors give. The present study's ten valid drought-recovery events are too few for such a comparison, which the authors acknowledge, and the comparison is the natural next test.

\section{Selection of Claim Type}

The recurrent gap invites an account of why it recurs. The following is a heuristic illustration of publishing pressures and not a formal economic derivation or an estimated model; the functional forms are there only to make visible which assumptions the argument needs. Fix a chain of increasingly demanding claims within a branch of Table~\ref{tab:ladder}, indexed $1,\dots,L$ by convention, and suppose the supported levels form an initial segment of the chain, with $e$ the highest of them. Let $\ell\in\{e,e+1,\dots,L\}$ be the level asserted in a given section $s$ of a paper (title, abstract, discussion, limitations). The author gains a section-specific reward $R_s(\ell)$, increasing in $\ell$, and bears a section-specific expected penalty $P_s(\ell-e)$ with $P_s(0)=0$, from reader and reviewer scrutiny. With discrete levels, asserting $k$ levels above $e$ is preferred to asserting none exactly when $R_s(e+k)-R_s(e)>P_s(k)$, and nothing in the model guarantees that the next available level is worthwhile. A continuous relaxation illustrates the local incentive. If $R_s$ is increasing and concave, $P_s$ is convex with $P_s'(0^+)=0$ (hedged overreach draws almost no marginal penalty), and $e<L$, then $U_s=R_s(\ell)-P_s(\ell-e)$ is concave with $U_s'(e)=R_s'(e)>0$, so the maximizer satisfies $\ell^\ast>e$.

The conclusion that $\ell^\ast>e$ is close to a restatement of the assumptions, and the model's use is to show which assumptions it needs, namely a positive marginal reward and a vanishing marginal penalty for small hedged excess. This is a statement about incentive structure, not motive. The claim that the excess is larger in titles and abstracts than in discussions requires an assumption absent from the displayed model, namely that $R_s$ is steeper and $P_s$ flatter in the title and abstract than in the discussion, since abstracts are read by more people and scrutinized less. If one assumes this, the model implies the following ordering: across a sample, the claim type coded from titles and abstracts exceeds that coded from results by more than the type coded from the discussion does. The ordering is a consequence of the assumption and would be testable only if someone coded such a sample; it is not a finding. Related work on exaggeration, spin, and selection for positive results gives context \cite{smaldino2016,vinkers2015,chiu2017,fraser2018}, though none of it tests this model. Without calibration against such a sample, the model shows only that a gap does not require any author to act in bad faith.

\section{What Would Count as Demonstrated Function}

For each claim beyond occupancy there is a test that could fail.

A nursery claim requires that juveniles be identified on the habitat, that their survival and growth there be estimated, and that they be linked to adult populations through otolith chemistry, tagging, or genetic assignment, in the sense of the contribution framework \cite{beck2001} and the concepts developed in the critical literature \cite{sheaves2006,dahlgren2006}. A connectivity claim requires movement or recruitment data in place of shared composition, and the treatment of structural and functional connectivity in \cite{taylor2006} is the relevant background.

For the distinctive peat-effect estimand, one suitable design for a claim of legacy-mediated function is a hard-substrate non-peat control matched on depth and exposure, with enough sites per level that the permutation floor and the intraclass correlation permit the effect of interest to be resolved; the definitions allow other counterfactuals and other credible identification designs, and a claim of functional inheritance additionally requires independent evidence about the earlier system's function and mechanism. For resilience capacity, the corresponding requirement is that capacity measures be validated against several disturbance outcomes across independent events, with the index evaluated on held-out events against specified alternative aggregations and component-based models, including forms that do not hard-code the ordering of importance.

None of these tests is exotic. Their cost is what distinguishes them from the characterizations that precede them.

\section{Limits of the Argument}

Two cases do not make a pattern. They are chosen as clear examples of a structure, and the argument shows what the structure is, not how common it is. The permutation-floor and design-effect arguments are easy to overapply. A limited design may still support a claim whose scope is restricted to the sampled system, and the pseudoreplication charge depends on the inferential target \cite{oksanen2001,davies2015}. The authors of both studies state limitations, and a reader should weigh those statements as part of the evidence. The claim gap is defined by a coding rule and could be disputed at the margins, and the ordering within Table~\ref{tab:ladder} is a convention and not a theorem. The selection model is uncalibrated. The characterization of the validation literature in particular rests on abstracts and excerpts and should be taken as a pointer to the sources and not as a summary that replaces them.

\section{Conclusion}

Both studies characterize candidate enabling conditions. The peat study shows that a consolidated substrate of terrestrial origin exists offshore and that organisms occupy it. The urban resilience study shows how a particular set of capacities can be combined into an index and how that index differed between two endpoints. The observed substrates and constructed capacities do not, by themselves, identify the proposed functions. Whether either study additionally demonstrates enabling effects, realized performance, or causal inheritance must be established claim by claim, and for that the tests of Section~8 are the relevant evidence. Constructed capacities, in particular, should not be treated as already validated possibilities. The practice that makes the mineral-evolution comparison checkable, stating the definition, the type of claim, and the evidence that would count, is available to the fields discussed here.

\paragraph{Note on sources.} Quoted statements from the two Hazen papers and the two case papers were checked against the full text. Bibliographic details of the cited works were checked against publisher, repository, or preprint records. The secondary literature on resilience validation was read through abstracts and excerpts.

\appendix
\section{Rounding bounds for the decomposition in Case B}
The reported values are $\overline{\mathrm{TE4}}$: $0.1184\to0.1168$; $\bar D$: $0.3723\to0.3424$; $\bar P$: $0.3149\to0.3441$. The terms are $\Delta(\bar D\bar P)=\bar D_1\bar P_1-\bar D_0\bar P_0=0.117820-0.117237=+0.000583$ and $\Delta r=(\overline{\mathrm{TE4}}_1-\bar D_1\bar P_1)-(\overline{\mathrm{TE4}}_0-\bar D_0\bar P_0)=-0.001020-0.001163=-0.002183$, which sum to $-0.0016$. Each reported value is rounded to four decimals, so the true value lies within $\pm0.00005$ of it. The six inputs are the six reported values (the TE4 mean, $\bar D$, and $\bar P$ at each endpoint). Each term is affine in each input separately, since $\Delta r$ is affine in the two TE4 values and bilinear in the $\bar D$ and $\bar P$ values, so its extrema over the rounding box occur at the box's $2^6=64$ corners. Enumerating the corners gives $\Delta(\bar D\bar P)\in[+0.00051,+0.00065]$ and $\Delta r\in[-0.00235,-0.00201]$, so the signs do not depend on rounding.

\begin{thebibliography}{99}
\bibitem{albertson2024} L.~K. Albertson, L.~S. Sklar, B.~B. Tumolo, W.~F. Cross, S.~F. Collins, H.~A. Woods, The ghosts of ecosystem engineers: legacy effects of biogenic modifications, \emph{Functional Ecology} 38, 52-72 (2024), doi:10.1111/1365-2435.14222.
\bibitem{bakkensen2017} L.~A. Bakkensen, C. Fox-Lent, L.~K. Read, I. Linkov, Validating resilience and vulnerability indices in the context of natural disasters, \emph{Risk Analysis} 37, 982-1004 (2017), doi:10.1111/risa.12677.
\bibitem{beck2001} M.~W. Beck et al., The identification, conservation, and management of estuarine and marine nurseries for fish and invertebrates, \emph{BioScience} 51, 633-641 (2001).
\bibitem{carpenter2001} S. Carpenter, B. Walker, J.~M. Anderies, N. Abel, From metaphor to measurement: resilience of what to what?, \emph{Ecosystems} 4, 765-781 (2001).
\bibitem{chiu2017} K. Chiu, Q. Grundy, L. Bero, ``Spin'' in published biomedical literature: a methodological systematic review, \emph{PLOS Biology} 15, e2002173 (2017), doi:10.1371/journal.pbio.2002173.
\bibitem{dahlgren2006} C.~P. Dahlgren et al., Marine nurseries and effective juvenile habitats: concepts and applications, \emph{Marine Ecology Progress Series} 312, 291-295 (2006).
\bibitem{davies2015} G.~M. Davies, A. Gray, Don't let spurious accusations of pseudoreplication limit our ability to learn from natural experiments (and other messy kinds of ecological monitoring), \emph{Ecology and Evolution} 5, 5295-5304 (2015), doi:10.1002/ece3.1782.
\bibitem{fraser2018} H. Fraser, T. Parker, S. Nakagawa, A. Barnett, F. Fidler, Questionable research practices in ecology and evolution, \emph{PLOS ONE} 13, e0200303 (2018), doi:10.1371/journal.pone.0200303.
\bibitem{gu2023} D. Gu, M. Dillard, M. Gerst, J. Loerzel, Validating commonly-used indicators for community resilience measurement, \emph{Natural Hazards Review} (2023), doi:10.1061/NHREFO.NHENG-1642.
\bibitem{gu2025} D. Gu, M. Dillard, M. Gerst, Assessment of community resilience frameworks: a validation through community outcomes after disaster events, NIST publication (listed as a talk), 10 December 2025.
\bibitem{halekotte2024} L. Halekotte, A. Mentges, D. Lichte, Do we practice what we preach? The dissonance between resilience understanding and measurement, preprint, arXiv:2410.07253 (2024).
\bibitem{hazen2008} R.~M. Hazen, D. Papineau, W. Bleeker, R.~T. Downs, J.~M. Ferry, T.~J. McCoy, D.~A. Sverjensky, H. Yang, Mineral evolution, \emph{American Mineralogist} 93, 1693-1720 (2008).
\bibitem{hazen2010} R.~M. Hazen, J.~M. Ferry, Mineral evolution: mineralogy in the fourth dimension, \emph{Elements} 6, 9-12 (2010), doi:10.2113/gselements.6.1.9.
\bibitem{hazen2022} R.~M. Hazen, S.~M. Morrison, On the paragenetic modes of minerals: a mineral evolution perspective, \emph{American Mineralogist} 107, 1262-1287 (2022).
\bibitem{hou2026} Y. Hou, Q. Hu, \"U. Halik, Structural diagnosis and evolution of urban resilience in dryland oasis urban agglomerations from an eco-immunological perspective, \emph{Scientific Reports} (2026), Article in Press, doi:10.1038/s41598-026-75595-w.
\bibitem{hurlbert1984} S.~H. Hurlbert, Pseudoreplication and the design of ecological field experiments, \emph{Ecological Monographs} 54, 187-211 (1984).
\bibitem{kish1965} L. Kish, \emph{Survey Sampling} (Wiley, New York, 1965).
\bibitem{laland2016} K.~N. Laland, B. Matthews, M.~W. Feldman, An introduction to niche construction theory, \emph{Evolutionary Ecology} 30, 191-202 (2016).
\bibitem{lebrato2026} M. Lebrato, J. Shannon, T. Russell, A. Jordan, M. Jije-Goncalves, M. da Gra\c{c}a, M. Caputo, A. Lillis, Ancient terrestrial forests reborn as modern marine ecosystems under shifting sea levels, \emph{Scientific Reports} 16, 24316 (2026), doi:10.1038/s41598-026-65430-7.
\bibitem{oksanen2001} L. Oksanen, Logic of experiments in ecology: is pseudoreplication a pseudoissue?, \emph{Oikos} 94, 27-38 (2001).
\bibitem{parnell2010} A.~C. Parnell, R. Inger, S. Bearhop, A.~L. Jackson, Source partitioning using stable isotopes: coping with too much variation, \emph{PLOS ONE} 5, e9672 (2010).
\bibitem{petersen2012} M.~L. Petersen, K.~E. Porter, S. Gruber, Y. Wang, M.~J. van der Laan, Diagnosing and responding to violations in the positivity assumption, \emph{Statistical Methods in Medical Research} 21, 31-54 (2012).
\bibitem{rubin1974} D.~B. Rubin, Estimating causal effects of treatments in randomized and nonrandomized studies, \emph{Journal of Educational Psychology} 66, 688-701 (1974).
\bibitem{sheaves2006} M. Sheaves, R. Baker, C.~R. Johnston, Marine nurseries and effective juvenile habitats: an alternative view, \emph{Marine Ecology Progress Series} 318, 303-306 (2006).
\bibitem{smaldino2016} P.~E. Smaldino, R. McElreath, The natural selection of bad science, \emph{Royal Society Open Science} 3, 160384 (2016), doi:10.1098/rsos.160384.
\bibitem{taylor2006} P.~D. Taylor, L. Fahrig, K.~A. With, Landscape connectivity: a return to the basics, in \emph{Connectivity Conservation}, K.~R. Crooks and M. Sanjayan, eds. (Cambridge University Press, 2006), pp. 29-43, doi:10.1017/CBO9780511754821.003.
\bibitem{vinkers2015} C.~H. Vinkers, J.~K. Tijdink, W.~M. Otte, Use of positive and negative words in scientific PubMed abstracts between 1974 and 2014: retrospective analysis, \emph{BMJ} 351, h6467 (2015).
\end{thebibliography}
\end{document}
